\item La rapidez del sonido en el aire (en metros por segundo) depende de la temperatura, de acuerdo con la expresión aproximada
$$ v = 331.5 + 0.607 T_C $$
donde $T_C$ es la temperatura Celsius. En aire seco, la temperatura disminuye casi 1 °C por cada 150 m de aumento en altura.
\begin{enumerate}
    \item Suponga que este cambio es constante hasta una altura de 9000 m. ¿Qué intervalo de tiempo se requiere para que el sonido de un avión que vuela a 9000 m llegue al suelo en un día cuando la temperatura del suelo es de 30 °C?
    \item ¿Qué pasaría si? Compare su respuesta con el intervalo de tiempo requerido si el aire estuviese uniformemente a 30 °C. ¿Qué intervalo de tiempo es más largo?
\end{enumerate}